% GENERATED FILE. Edit canonical lessons, then run npm run generate:books.
% canonical-source-hash: fnv1a64:14ab2779076be2a0
% canonical-lessons: HI-C233-uttar, HI-C233-daksin, HI-C233-purv, HI-C233-pascim, HI-C233-daya, HI-R233-first-pass-words-from-chapters-225-228, HI-R233-second-pass-words-from-chapters-229-233

\chapter{North, South, East, West, Right}
\label{ch:hi-north-south-east-west-right}
\hlchaptermodality{233}

\begin{chapteropening}
\textbf{By the end of this chapter:} \emph{I can say north, south, east, west and right.}

Complete the last lesson of chapter 233: I can say north, south, east, west and right.
\end{chapteropening}

\section[\texorpdfstring{\hi{उत्तर} (uttar)}{उत्तर (uttar)}]{\hi{उत्तर} (uttar) — north}

\label{lesson:HI-C233-uttar}



\begin{quote}
Before the new one: say the Hindi for to whisper, then the Hindi for to scold.
\end{quote}

\subsection*{You'll want to know: \hi{उत्तर}}
\textbf{\hi{उत्तर}} — \emph{uttar} — \textquotedblleft{}north\textquotedblright{}.

\begin{grammarlens}[title={What you've built}]
One more word for this chapter.
\end{grammarlens}

\subsection*{Your turn}
\begin{itemize}
\raggedright
  \item \emph{Say it:} \emph{uttar}
  \item \emph{Say it:} \emph{uttar}, once more
  \item \emph{Say it:} say \emph{uttar}
  \item \emph{Recall:} say the Hindi for to whisper, then the Hindi for to scold, then say \emph{uttar} again
\end{itemize}

\begin{tcolorbox}[breakable,colback=teal!4,colframe=teal!35!black,title={Before you move on}]
What does \hi{उत्तर} mean? (\textquotedblleft{}North\textquotedblright{}.) Say it once more.
\end{tcolorbox}
\section[\texorpdfstring{\hi{दक्षिण} (dakṣiṇ)}{दक्षिण (dakṣiṇ)}]{\hi{दक्षिण} (dakṣiṇ) — south}

\label{lesson:HI-C233-daksin}



\begin{quote}
Before the new one: say the Hindi for to scold, then the Hindi for north.
\end{quote}

\subsection*{You'll want to know: \hi{दक्षिण}}
\textbf{\hi{दक्षिण}} — \emph{dakṣiṇ} — \textquotedblleft{}south\textquotedblright{}.

\begin{grammarlens}[title={What you've built}]
One more word for this chapter.
\end{grammarlens}

\subsection*{Your turn}
\begin{itemize}
\raggedright
  \item \emph{Say it:} \emph{dakṣiṇ}
  \item \emph{Say it:} \emph{dakṣiṇ}, once more
  \item \emph{Say it:} say \emph{dakṣiṇ}
  \item \emph{Recall:} say the Hindi for to scold, then the Hindi for north, then say \emph{dakṣiṇ} again
\end{itemize}

\begin{tcolorbox}[breakable,colback=teal!4,colframe=teal!35!black,title={Before you move on}]
What does \hi{दक्षिण} mean? (\textquotedblleft{}South\textquotedblright{}.) Say it once more.
\end{tcolorbox}
\section[\texorpdfstring{\hi{पूर्व} (pūrv)}{पूर्व (pūrv)}]{\hi{पूर्व} (pūrv) — east}

\label{lesson:HI-C233-purv}



\begin{quote}
Before the new one: say the Hindi for north, then the Hindi for south.
\end{quote}

\subsection*{You'll want to know: \hi{पूर्व}}
\textbf{\hi{पूर्व}} — \emph{pūrv} — \textquotedblleft{}east\textquotedblright{}.

\begin{grammarlens}[title={What you've built}]
One more word for this chapter.
\end{grammarlens}

\subsection*{Your turn}
\begin{itemize}
\raggedright
  \item \emph{Say it:} \emph{pūrv}
  \item \emph{Say it:} \emph{pūrv}, once more
  \item \emph{Say it:} say \emph{pūrv}
  \item \emph{Recall:} say the Hindi for north, then the Hindi for south, then say \emph{pūrv} again
\end{itemize}

\begin{tcolorbox}[breakable,colback=teal!4,colframe=teal!35!black,title={Before you move on}]
What does \hi{पूर्व} mean? (\textquotedblleft{}East\textquotedblright{}.) Say it once more.
\end{tcolorbox}
\section[\texorpdfstring{\hi{पश्चिम} (paścim)}{पश्चिम (paścim)}]{\hi{पश्चिम} (paścim) — west}

\label{lesson:HI-C233-pascim}



\begin{quote}
Before the new one: say the Hindi for south, then the Hindi for east.
\end{quote}

\subsection*{You'll want to know: \hi{पश्चिम}}
\textbf{\hi{पश्चिम}} — \emph{paścim} — \textquotedblleft{}west\textquotedblright{}.

\begin{grammarlens}[title={What you've built}]
One more word for this chapter.
\end{grammarlens}

\subsection*{Your turn}
\begin{itemize}
\raggedright
  \item \emph{Say it:} \emph{paścim}
  \item \emph{Say it:} \emph{paścim}, once more
  \item \emph{Say it:} say \emph{paścim}
  \item \emph{Recall:} say the Hindi for south, then the Hindi for east, then say \emph{paścim} again
\end{itemize}

\begin{tcolorbox}[breakable,colback=teal!4,colframe=teal!35!black,title={Before you move on}]
What does \hi{पश्चिम} mean? (\textquotedblleft{}West\textquotedblright{}.) Say it once more.
\end{tcolorbox}
\section[\texorpdfstring{\hi{दायाँ} (dāyā̃)}{दायाँ (dāyā̃)}]{\hi{दायाँ} (dāyā̃) — right (side)}

\label{lesson:HI-C233-daya}



\begin{quote}
Before the new one: say the Hindi for east, then the Hindi for west.
\end{quote}

\subsection*{You'll want to know: \hi{दायाँ}}
\textbf{\hi{दायाँ}} — \emph{dāyā̃} — \textquotedblleft{}right (side)\textquotedblright{}.

\begin{grammarlens}[title={What you've built}]
One more word for this chapter.
\end{grammarlens}

\subsection*{Your turn}
\begin{itemize}
\raggedright
  \item \emph{Say it:} \emph{dāyā̃}
  \item \emph{Say it:} \emph{dāyā̃}, once more
  \item \emph{Say it:} say \emph{dāyā̃}
  \item \emph{Recall:} say the Hindi for east, then the Hindi for west, then say \emph{dāyā̃} again
\end{itemize}

\begin{tcolorbox}[breakable,colback=teal!4,colframe=teal!35!black,title={Before you move on}]
What does \hi{दायाँ} mean? (\textquotedblleft{}Right (side)\textquotedblright{}.) Say it once more.
\end{tcolorbox}
\section[(dialogue)]{Words from chapters 225-228 — a first pass}

\label{lesson:HI-R233-first-pass-words-from-chapters-225-228}



\begin{quote}
Say the words for west and right.
\end{quote}

\subsection*{Your turn}
Say these aloud:

\begin{itemize}
\raggedright
  \item to clean, to dust, to wipe, to wring, to dry
  \item to dry, to go wrong, to make, to be made, to break
  \item to tear, to tear, to bring up, to split, to sow
  \item to grow, to grow, to hand over, to press, to turn
\end{itemize}

\begin{tcolorbox}[breakable,colback=teal!4,colframe=teal!35!black,title={Before you move on}]
To clean? (\textbf{\hi{साफ़} \hi{करना}}, \emph{sāf karnā}.) To dust? (\textbf{\hi{झाड़ना}}, \emph{jhāṛnā}.) To wipe? (\textbf{\hi{पोंछना}}, \emph{põchnā}.) To wring? (\textbf{\hi{निचोड़ना}}, \emph{nicoṛnā}.) To dry? (\textbf{\hi{सुखाना}}, \emph{sukhānā}.) To dry? (\textbf{\hi{सूखना}}, \emph{sūkhnā}.) To go wrong? (\textbf{\hi{बिगड़ना}}, \emph{bigaṛnā}.) To make? (\textbf{\hi{बनाना}}, \emph{banānā}.) To be made? (\textbf{\hi{बनना}}, \emph{bannā}.) To break? (\textbf{\hi{टूटना}}, \emph{ṭūṭnā}.) To tear? (\textbf{\hi{फटना}}, \emph{phaṭnā}.) To tear? (\textbf{\hi{फाड़ना}}, \emph{phāṛnā}.) To bring up? (\textbf{\hi{पालना}}, \emph{pālnā}.) To split? (\textbf{\hi{चीरना}}, \emph{cīrnā}.) To sow? (\textbf{\hi{बोना}}, \emph{bonā}.) To grow? (\textbf{\hi{उगना}}, \emph{ugnā}.) To grow? (\textbf{\hi{उगाना}}, \emph{ugānā}.) To hand over? (\textbf{\hi{पकड़ाना}}, \emph{pakṛānā}.) To press? (\textbf{\hi{दबाना}}, \emph{dabānā}.) To turn? (\textbf{\hi{घुमाना}}, \emph{ghumānā}.)
\end{tcolorbox}
\section[(dialogue)]{Words from chapters 229-233 — a second pass}

\label{lesson:HI-R233-second-pass-words-from-chapters-229-233}



\begin{quote}
Say the words for to bend and to gather up.
\end{quote}

\subsection*{Your turn}
Say these aloud:

\begin{itemize}
\raggedright
  \item to bend, to gather up, to spread, to spread, to be left behind
  \item to bend down, to hang, to hang, to float, to sink
  \item to flow, to drip, to sneeze, to yawn, to cough
  \item to shiver, to smile, to scream, to whisper, to scold
  \item north, south, east, west, right
\end{itemize}

\begin{tcolorbox}[breakable,colback=teal!4,colframe=teal!35!black,title={Before you move on}]
To bend? (\textbf{\hi{मोड़ना}}, \emph{moṛnā}.) To gather up? (\textbf{\hi{समेटना}}, \emph{sameṭnā}.) To spread? (\textbf{\hi{फैलाना}}, \emph{phailānā}.) To spread? (\textbf{\hi{फैलना}}, \emph{phailnā}.) To be left behind? (\textbf{\hi{छूटना}}, \emph{chūṭnā}.) To bend down? (\textbf{\hi{झुकना}}, \emph{jhuknā}.) To hang? (\textbf{\hi{लटकना}}, \emph{laṭaknā}.) To hang? (\textbf{\hi{लटकाना}}, \emph{laṭkānā}.) To float? (\textbf{\hi{तैराना}}, \emph{tairānā}.) To sink? (\textbf{\hi{डूबना}}, \emph{ḍūbnā}.) To flow? (\textbf{\hi{बहना}}, \emph{bahnā}.) To drip? (\textbf{\hi{टपकना}}, \emph{ṭapaknā}.) To sneeze? (\textbf{\hi{छींकना}}, \emph{chīṅknā}.) To yawn? (\textbf{\hi{जम्हाई} \hi{लेना}}, \emph{jamhāī lenā}.) To cough? (\textbf{\hi{खाँसना}}, \emph{khā̃snā}.) To shiver? (\textbf{\hi{काँपना}}, \emph{kā̃pnā}.) To smile? (\textbf{\hi{मुस्कुराना}}, \emph{muskurānā}.) To scream? (\textbf{\hi{चीखना}}, \emph{cīkhnā}.) To whisper? (\textbf{\hi{फुसफुसाना}}, \emph{phusphusānā}.) To scold? (\textbf{\hi{डाँटना}}, \emph{ḍā̃ṭnā}.) North? (\textbf{\hi{उत्तर}}, \emph{uttar}.) South? (\textbf{\hi{दक्षिण}}, \emph{dakṣiṇ}.) East? (\textbf{\hi{पूर्व}}, \emph{pūrv}.) West? (\textbf{\hi{पश्चिम}}, \emph{paścim}.) Right? (\textbf{\hi{दायाँ}}, \emph{dāyā̃}.)
\end{tcolorbox}
